% Chapter: Sheaf Types and Cohomological Descent
% Part of the CCTT Monograph — Proof Calculus, Chapter III

\chapter{Sheaf Types and Cohomological Descent}
\label{ch:c4-descent}

Chapter~\ref{ch:c4-calculus} introduced the descent rule as a typing rule.
This chapter develops the full sheaf-theoretic infrastructure: the site
category for CCTT, presheaf and sheaf types, the \v{C}ech nerve
construction \emph{within} $\mathrm{C}^4$, $\Hone$ as a type, and
the descent theorem with a complete proof.  We prove that
the CCTT equivalence checker's per-fiber Z3 + \v{C}ech $\Hone$ strategy is
a special case of the descent rule, establishing the soundness of the
implementation with respect to the calculus.


% ─────────────────────────────────────────────────────────
\section{The Site Category for CCTT}
\label{sec:site-category}

\begin{definition}[The CCTT Site $\site_\mathrm{CCTT}$]
\label{def:cctt-site}
The \textbf{CCTT site} $\site_\mathrm{CCTT}$ is a category equipped with a
Grothendieck topology:
\begin{enumerate}
\item \textbf{Objects:} duck-type fibers $\tau \in \{\mathsf{TInt},
  \mathsf{TFloat}, \mathsf{TStr}, \mathsf{TBool}, \mathsf{TList},
  \mathsf{TDict}, \mathsf{TSet}, \mathsf{TTuple}, \mathsf{TNone},
  \mathsf{TCallable}, \mathsf{TAny}\}$.

\item \textbf{Morphisms:} subtyping inclusions $\iota_{\tau \to \tau'}$
  where $\tau$ is a subtype of $\tau'$.  The morphisms form a partial
  order with $\mathsf{TAny}$ as the top element:
  \[
    \mathsf{TInt} \hookrightarrow \mathsf{TFloat}
    \hookrightarrow \mathsf{TAny}, \qquad
    \mathsf{TBool} \hookrightarrow \mathsf{TInt}, \qquad
    \mathsf{TList} \hookrightarrow \mathsf{TAny}, \quad \ldots
  \]

\item \textbf{Covering families:} A family $\{U_\alpha \to U\}$ is a
  cover if and only if $U = \bigcup_\alpha U_\alpha$ (every element of
  $U$ factors through some $U_\alpha$).  For the terminal object
  $\mathsf{TAny}$, the \emph{standard cover} is
  $\covers_\mathrm{std} = \{\mathsf{TInt}, \mathsf{TFloat},
  \mathsf{TStr}, \mathsf{TBool}, \mathsf{TList}, \mathsf{TDict},
  \mathsf{TSet}, \mathsf{TTuple}, \mathsf{TNone}, \mathsf{TCallable}\}$.
\end{enumerate}
\end{definition}

\begin{remark}
The site $\site_\mathrm{CCTT}$ is \emph{finite}---it has finitely many
objects and morphisms.  This is crucial for decidability of type checking
(Chapter~\ref{ch:c4-metatheory}): \v{C}ech cohomology over a finite site
is computable.
\end{remark}


% ─────────────────────────────────────────────────────────
\section{Presheaf Types}
\label{sec:presheaf-types}

\begin{definition}[Presheaf Type]
\label{def:presheaf-type}
A \textbf{presheaf type} over $\site$ is a contravariant functor
$\presheaf : \site^{\mathrm{op}} \to \mathcal{U}_i$.  Concretely, it
assigns:
\begin{enumerate}
\item To each object $U \in \site$, a type $\presheaf(U) : \mathcal{U}_i$.
\item To each morphism $f : V \to U$ in $\site$, a restriction map
  $\restr{}{f} : \presheaf(U) \to \presheaf(V)$.
\item These satisfy \textbf{functoriality}:
  $\restr{}{\mathrm{id}} = \mathrm{id}$ and
  $\restr{}{g \circ f} = \restr{}{f} \circ \restr{}{g}$.
\end{enumerate}
\end{definition}

\begin{example}[OTerm Behavior as a Presheaf]
Given an OTerm $t \in \mathcal{U}_\PyObj$, its \textbf{behavioral
presheaf} $\mathcal{B}_t$ assigns to each fiber $\tau$:
\[
  \mathcal{B}_t(\tau) \defeq \{\text{input-output pairs of $t$ on
    inputs of type $\tau$}\}
\]
The restriction along $\iota_{\tau \to \tau'} : \tau \hookrightarrow
\tau'$ simply restricts the domain: $\restr{\mathcal{B}_t(\tau')}{\tau}
\subseteq \mathcal{B}_t(\tau)$ (a function that works on $\tau'$ inputs
also works on $\tau$ inputs by subtyping).
\end{example}


% ─────────────────────────────────────────────────────────
\section{The Sheaf Condition and Sheaf Types}
\label{sec:sheaf-types}

\begin{definition}[Sheaf Condition]
\label{def:c4-sheaf-condition}
A presheaf $\presheaf$ over $\site$ is a \textbf{sheaf} if for every
covering family $\{U_\alpha \to U\}_{\alpha \in I}$ and every
\emph{compatible family} of local sections
$s_\alpha \in \presheaf(U_\alpha)$ satisfying
\[
  \restr{s_\alpha}{U_\alpha \cap U_\beta} =
  \restr{s_\beta}{U_\alpha \cap U_\beta}
  \qquad \text{for all } \alpha, \beta,
\]
there exists a \emph{unique} global section $s \in \presheaf(U)$ with
$\restr{s}{U_\alpha} = s_\alpha$ for all $\alpha$.
\end{definition}

\begin{definition}[Sheaf Type in $\mathrm{C}^4$]
\label{def:sheaf-type-c4}
A type $A : \mathcal{U}_i^{\site}$ is a \textbf{sheaf type} if it
satisfies the sheaf condition internally: the canonical map
\[
  A(U) \xrightarrow{\;\prod_\alpha \restr{}{U_\alpha}\;}
  \prod_\alpha A(U_\alpha)
  \;\rightrightarrows\;
  \prod_{\alpha,\beta} A(U_\alpha \cap U_\beta)
\]
is an equalizer in $\mathcal{U}_i$ for every cover $\{U_\alpha\}$ of $U$.
In $\mathrm{C}^4$, the sheaf condition is expressed as a propositional
identity type: a term witnessing that the equalizer diagram is exact.
\end{definition}


% ─────────────────────────────────────────────────────────
\section{The \v{C}ech Nerve Within $\mathrm{C}^4$}
\label{sec:cech-nerve}

\begin{definition}[\v{C}ech Nerve]
\label{def:cech-nerve}
Given a cover $\covers = \{U_\alpha\}_{\alpha \in I}$ of $U$, the
\textbf{\v{C}ech nerve} is the simplicial object:
\begin{align*}
  C_0(\covers) &\defeq \coprod_\alpha U_\alpha \\
  C_1(\covers) &\defeq \coprod_{\alpha,\beta}
    (U_\alpha \cap U_\beta) \\
  C_2(\covers) &\defeq \coprod_{\alpha,\beta,\gamma}
    (U_\alpha \cap U_\beta \cap U_\gamma) \\
  &\;\;\vdots
\end{align*}
with face maps $d_i$ given by the inclusion morphisms and degeneracy maps
$s_i$ by the diagonal.
\end{definition}

\begin{definition}[\v{C}ech Complex]
\label{def:cech-complex}
For a presheaf $\presheaf$ on $\site$ and a cover $\covers$, the
\textbf{\v{C}ech complex} $\check{C}^\bullet(\covers, \presheaf)$ is:
\[
  \prod_\alpha \presheaf(U_\alpha)
  \xrightarrow{\;\delta^0\;}
  \prod_{\alpha,\beta} \presheaf(U_\alpha \cap U_\beta)
  \xrightarrow{\;\delta^1\;}
  \prod_{\alpha,\beta,\gamma}
    \presheaf(U_\alpha \cap U_\beta \cap U_\gamma)
  \xrightarrow{\;\delta^2\;} \cdots
\]
where $\delta^0(s)_{\alpha\beta} = \restr{s_\beta}{U_{\alpha\beta}}
- \restr{s_\alpha}{U_{\alpha\beta}}$ and $\delta^1$ is the alternating
sum on triple overlaps.  The \textbf{\v{C}ech cohomology} is
$\check{H}^n(\covers, \presheaf) \defeq \ker \delta^n / \mathrm{im}\,
\delta^{n-1}$.
\end{definition}

\begin{definition}[$\Hone$ as a Type]
\label{def:h1-type}
In $\mathrm{C}^4$, we internalize $\Hone$ as a type:
\[
  \Hone(\covers, \presheaf) \;\defeq\;
  \frac{\Sigma(g : \prod_{\alpha\beta}
    \presheaf(U_{\alpha\beta})).\; \delta^1(g) = 0}
  {\sim_{\mathrm{coboundary}}}
\]
where $g \sim_\mathrm{coboundary} g'$ iff there exists
$h : \prod_\alpha \presheaf(U_\alpha)$ with $g'_{\alpha\beta} -
g_{\alpha\beta} = \restr{h_\beta}{U_{\alpha\beta}} -
\restr{h_\alpha}{U_{\alpha\beta}}$.

A proof that $\Hone(\covers, \presheaf)$ is contractible (has a unique
element, namely $0$) is a witness that all cocycles are coboundaries.
\end{definition}


% ─────────────────────────────────────────────────────────
\section{The Descent Theorem}
\label{sec:descent-theorem}

We now prove the descent theorem in full detail within $\mathrm{C}^4$.

\begin{theorem}[Descent]
\label{thm:c4-descent-full}
Let $A : \mathcal{U}_i^{\site}$ be a sheaf type, $\covers =
\{U_\alpha\}$ a covering of the terminal object.  The following are
equivalent:
\begin{enumerate}
\item \textbf{(Descent data glues.)} Every compatible family
  $(t_\alpha)_\alpha$ with cocycle $(p_{\alpha\beta})$ satisfying
  the cocycle condition extends to a global section $t : A$.
\item \textbf{($\Hone$ vanishes.)}
  $\Hone(\covers, \underline{\mathrm{Aut}}(A)) = 0$, where
  $\underline{\mathrm{Aut}}(A)$ is the sheaf of automorphisms of $A$.
\end{enumerate}
\end{theorem}

\begin{proof}
\textbf{$(2) \Rightarrow (1)$:}
Suppose $\Hone(\covers, \underline{\mathrm{Aut}}(A)) = 0$.  Given
descent data $(t_\alpha, p_{\alpha\beta})$ where:
\begin{itemize}
\item $t_\alpha : \restr{A}{U_\alpha}$ for each $\alpha$,
\item $p_{\alpha\beta} :
  \restr{t_\alpha}{U_{\alpha\beta}} \equiv
  \restr{t_\beta}{U_{\alpha\beta}}$ for each $\alpha, \beta$,
\item The cocycle condition:
  $\restr{p_{\alpha\gamma}}{U_{\alpha\beta\gamma}} =
  \restr{p_{\beta\gamma}}{U_{\alpha\beta\gamma}} \circ
  \restr{p_{\alpha\beta}}{U_{\alpha\beta\gamma}}$.
\end{itemize}

We interpret the compatibility witnesses $p_{\alpha\beta}$ as paths.
On each double overlap $U_{\alpha\beta}$, the path $p_{\alpha\beta}$
gives a ``transition map'' between the local sections.  The cocycle
condition says these transition maps compose correctly on triple overlaps.

Define the 1-cocycle $g_{\alpha\beta} \in
\mathrm{Aut}(\restr{A}{U_{\alpha\beta}})$ by
$g_{\alpha\beta}(x) = p_{\alpha\beta}^{-1} \circ x \circ
p_{\alpha\beta}$, where we use path transport.  The cocycle condition on
$p$ gives the cocycle condition $g_{\alpha\gamma} =
g_{\beta\gamma} \circ g_{\alpha\beta}$ on triple overlaps.

Since $\Hone(\covers, \underline{\mathrm{Aut}}(A)) = 0$, this cocycle
is a coboundary: there exist $h_\alpha \in
\mathrm{Aut}(\restr{A}{U_\alpha})$ with
$g_{\alpha\beta} = h_\beta^{-1} \circ h_\alpha$ on $U_{\alpha\beta}$.

Define $\tilde{t}_\alpha \defeq h_\alpha(t_\alpha) :
\restr{A}{U_\alpha}$.  Then on $U_{\alpha\beta}$:
\begin{align*}
  \restr{\tilde{t}_\alpha}{U_{\alpha\beta}}
  &= \restr{(h_\alpha(t_\alpha))}{U_{\alpha\beta}} \\
  &= h_\alpha(\restr{t_\alpha}{U_{\alpha\beta}}) \\
  &= h_\alpha(p_{\alpha\beta}(\restr{t_\beta}{U_{\alpha\beta}}))
  \tag{by compatibility} \\
  &= (h_\alpha \circ g_{\alpha\beta}^{-1})(
    \restr{t_\beta}{U_{\alpha\beta}}) \\
  &= h_\beta(\restr{t_\beta}{U_{\alpha\beta}})
  \tag{since $g_{\alpha\beta} = h_\beta^{-1} h_\alpha$} \\
  &= \restr{\tilde{t}_\beta}{U_{\alpha\beta}}
\end{align*}

So $\tilde{t}_\alpha$ and $\tilde{t}_\beta$ agree on every overlap, and
since $A$ is a sheaf, the sections $(\tilde{t}_\alpha)$ glue to a unique
global section $t : A$.

\textbf{$(1) \Rightarrow (2)$:}
Suppose every descent datum glues.  Let $(g_{\alpha\beta})$ be a 1-cocycle
in $\underline{\mathrm{Aut}}(A)$.  Pick any global section $t_0 : A$
(if $A$ has a global section; otherwise, pick local sections on a
refinement).  Define $t_\alpha \defeq \restr{t_0}{U_\alpha}$ and
$p_{\alpha\beta} \defeq g_{\alpha\beta}(\restr{t_0}{U_{\alpha\beta}})$.
By the cocycle condition on $g$, the $p_{\alpha\beta}$ satisfy the
cocycle condition for descent data.  By hypothesis~(1), the modified
sections $g_{\alpha\beta}(t_\alpha)$ glue to a global section $t'$.
Setting $h_\alpha \defeq t' \circ t_\alpha^{-1}$ on $U_\alpha$, we
get $g_{\alpha\beta} = h_\beta^{-1} \circ h_\alpha$, so $g$ is a
coboundary.  Hence $\Hone = 0$.
\end{proof}


% ─────────────────────────────────────────────────────────
\section{Mayer--Vietoris in $\mathrm{C}^4$}
\label{sec:mayer-vietoris}

\begin{theorem}[Mayer--Vietoris Exact Sequence]
\label{thm:mayer-vietoris}
For a sheaf $\presheaf$ over $\site$ and a cover $\covers =
\{U_1, U_2\}$ with two elements, there is a long exact sequence:
\[
  0 \to \Hzero(\covers, \presheaf) \to
  \presheaf(U_1) \oplus \presheaf(U_2) \to
  \presheaf(U_1 \cap U_2) \xrightarrow{\;\delta\;}
  \Hone(\covers, \presheaf) \to 0
\]
In $\mathrm{C}^4$, this is a sequence of types and functions with
the property that the image of each map equals the kernel of the next.
\end{theorem}

\begin{proof}
The \v{C}ech complex for a two-element cover is:
\[
  \presheaf(U_1) \times \presheaf(U_2) \xrightarrow{\delta^0}
  \presheaf(U_1 \cap U_2) \xrightarrow{\delta^1} 0
\]
since there are no triple overlaps (the cover has only two elements).
The coboundary map is $\delta^0(s_1, s_2) = \restr{s_2}{U_{12}} -
\restr{s_1}{U_{12}}$.

Then:
\begin{itemize}
\item $\Hzero = \ker \delta^0 = \{(s_1, s_2) \mid
  \restr{s_1}{U_{12}} = \restr{s_2}{U_{12}}\}$ = global sections of
  $\presheaf$.
\item $\Hone = \presheaf(U_{12}) / \mathrm{im}(\delta^0)$ = the
  cokernel, measuring how many ``local data on the overlap'' fail to
  extend.
\end{itemize}

The exact sequence
$0 \to \Hzero \hookrightarrow \presheaf(U_1) \oplus \presheaf(U_2)
\xrightarrow{\delta^0} \presheaf(U_{12}) \twoheadrightarrow \Hone \to 0$
is immediate from the definitions.
\end{proof}

\begin{example}[Mayer--Vietoris for Integer/String Fibers]
Consider the CCTT standard cover restricted to two fibers:
$U_1 = \mathsf{TInt}$, $U_2 = \mathsf{TStr}$.  For a function $f$,
the presheaf of behaviors gives:
\begin{itemize}
\item $\mathcal{B}_f(\mathsf{TInt})$ = behavior on integer inputs
\item $\mathcal{B}_f(\mathsf{TStr})$ = behavior on string inputs
\item $\mathcal{B}_f(\mathsf{TInt} \cap \mathsf{TStr}) = \emptyset$
  (no overlap between integers and strings)
\end{itemize}
Since $\presheaf(U_{12}) = 0$, we get $\Hone = 0$, and
$\Hzero = \presheaf(U_1) \oplus \presheaf(U_2)$.  In other words:
disjoint fibers always glue.  This is a formal justification of the
``decompose and conquer'' strategy.
\end{example}


% ─────────────────────────────────────────────────────────
\section{Obstruction Theory: $\Hone \neq 0$ as Type Error}
\label{sec:obstruction}

\begin{definition}[Obstruction Class]
\label{def:obstruction}
When the descent rule \emph{fails}---when local proofs on fibers
cannot be glued into a global proof---the \v{C}ech cohomology class
$[g] \in \Hone(\covers, \underline{\mathrm{Aut}}(A))$ is the
\textbf{obstruction class}.  It records \emph{which} overlaps are
inconsistent and \emph{how}.
\end{definition}

\begin{theorem}[Obstruction as Diagnostic]
\label{thm:obstruction-diagnostic}
Given an attempted descent with local sections $(t_\alpha)$ and
compatibility witnesses $(p_{\alpha\beta})$ that fail the cocycle
condition, the non-trivial class $[g] \in \Hone$ decomposes as:
\[
  [g] = \sum_{\alpha < \beta < \gamma}
  \underbrace{[\restr{p_{\alpha\gamma}}{U_{\alpha\beta\gamma}}
  - \restr{p_{\beta\gamma}}{U_{\alpha\beta\gamma}} \circ
  \restr{p_{\alpha\beta}}{U_{\alpha\beta\gamma}}]}_{\text{cocycle
  failure on triple overlap $(\alpha,\beta,\gamma)$}}
\]
Each summand pinpoints a specific triple of fibers where gluing fails.
\end{theorem}

\begin{proof}
By definition, the cocycle condition requires $\delta^1(p) = 0$ where
$\delta^1$ is the \v{C}ech coboundary on 1-cochains.  If $\delta^1(p)
\neq 0$, then $\delta^1(p)$ is a 2-cochain whose components on each
triple overlap $U_{\alpha\beta\gamma}$ measure the failure.  The
class $[g] = [\delta^1(p)] \in \Htwo$ would be the obstruction, but
for a finite site with few overlaps, the relevant obstruction lives
in $\Hone$ (since the non-abelian case reduces to $\Hone$ for the
automorphism sheaf).  In the abelian approximation, the components are
as stated.
\end{proof}

\begin{example}[Obstruction in Practice]
Consider proving that functions $f$ and $g$ are equivalent.  Suppose
the per-fiber proofs succeed on $\mathsf{TInt}$ and $\mathsf{TList}$
but the overlap $\mathsf{TList}[\mathsf{TInt}]$ (lists of integers)
has an inconsistency: the proof on $\mathsf{TInt}$ expects
$f(\texttt{[1,2,3]}) = g(\texttt{[1,2,3]})$ via path $p_1$, while the
proof on $\mathsf{TList}$ expects it via path $p_2 \neq p_1$.

The obstruction class reports:
\begin{quote}
  \textsf{Gluing failed on overlap TList[TInt]:}\\
  \textsf{\ \ fiber TInt proof uses axiom D3 (associativity)}\\
  \textsf{\ \ fiber TList proof uses axiom D17 (map-fold fusion)}\\
  \textsf{\ \ these give different normal forms on the overlap.}
\end{quote}

This diagnostic is qualitatively richer than ``proof failed'' or
``tactic timeout.''
\end{example}


% ─────────────────────────────────────────────────────────
\section{Soundness of the Cohomological Equivalence Checker}
\label{sec:checker-soundness}

We now prove that the CCTT equivalence checker (as implemented in
\texttt{new\_src/cctt/proof\_theory/}) is sound with respect to the
$\mathrm{C}^4$ descent rule.

\begin{theorem}[Checker Soundness]
\label{thm:checker-soundness}
The CCTT equivalence checker's strategy:
\begin{enumerate}
\item Decompose the proof obligation into per-fiber obligations via the
  duck-type fiber decomposition.
\item Prove each per-fiber obligation using Z3 (for structural properties)
  or axiom chains (for semantic properties).
\item Compute $\Hone$ of the \v{C}ech complex; if $\Hone = 0$, declare
  the proofs globally consistent.
\end{enumerate}
is a special case of the descent rule in $\mathrm{C}^4$.
\end{theorem}

\begin{proof}
We show that each step of the checker corresponds to a $\mathrm{C}^4$
derivation.

\textbf{Step 1: Fiber decomposition.}
The checker computes duck-type fibers $\tau_1, \ldots, \tau_k$ for the
input programs.  This corresponds to choosing the standard cover
$\covers_\mathrm{std}$ of $\mathsf{TAny}$ in $\site_\mathrm{CCTT}$
(Definition~\ref{def:cctt-site}) and applying the restriction rule (Res)
from Chapter~\ref{ch:c4-calculus} to obtain fiber-restricted obligations
$\Gamma \vdash_{\tau_j} \restr{\den{f}}{\tau_j} \equiv
\restr{\den{g}}{\tau_j}$.

\textbf{Step 2: Per-fiber proofs.}
For each fiber $\tau_j$:
\begin{itemize}
\item If Z3 proves $\restr{\den{f}}{\tau_j} = \restr{\den{g}}{\tau_j}$,
  this is a $\mathsf{Z3Discharge}$ proof term
  (Definition~\ref{def:z3-discharge} in Chapter~\ref{ch:proof-theory-foundations}),
  which corresponds to a path
  $p_j : \Path_{\restr{\mathcal{U}_\PyObj}{\tau_j}}\,
  \restr{\den{f}}{\tau_j}\, \restr{\den{g}}{\tau_j}$.
\item If an axiom chain $D_{i_1} \cdots D_{i_n}$ proves it, this
  corresponds to a HIT path composition by
  Theorem~\ref{thm:path-axiom}, restricted to fiber $\tau_j$.
\end{itemize}

\textbf{Step 3: \v{C}ech gluing.}
The checker computes $\Hone(\covers, \mathcal{B})$ where $\mathcal{B}$
is the behavioral presheaf.  If $\Hone = 0$, the per-fiber paths
$(p_j)$ satisfy the cocycle condition (the overlap conditions hold
trivially because Z3 proofs on overlaps agree by confluence of the
Z3 decision procedure).  The descent rule then produces
$\mathsf{desc}(\{p_j\}, \{\mathsf{refl}\}) :
\Path_{\mathcal{U}_\PyObj}\, \den{f}\, \den{g}$.

Hence every successful checker run produces a valid $\mathrm{C}^4$
typing derivation.
\end{proof}


% ─────────────────────────────────────────────────────────
\section{Advanced: Descent for Dependent Types}
\label{sec:descent-dependent}

\begin{definition}[Dependent Descent]
\label{def:dependent-descent}
For a type family $B : A \to \mathcal{U}_i^{\site}$ and local sections
$b_\alpha : \Pi(x : \restr{A}{U_\alpha}).\, \restr{B(x)}{U_\alpha}$
with compatibility witnesses
$q_{\alpha\beta} : \restr{b_\alpha}{U_{\alpha\beta}} \equiv
\restr{b_\beta}{U_{\alpha\beta}}$ (over the transport along $p_{\alpha\beta}$),
the \textbf{dependent descent} produces:
\[
  \mathsf{desc}^{\Sigma}(\{(t_\alpha, b_\alpha)\},
  \{(p_{\alpha\beta}, q_{\alpha\beta})\}) :
  \Sigma(x : A).\, B(x)
\]
provided the combined cocycle condition holds.
\end{definition}

\begin{proposition}
Dependent descent is reducible to ordinary descent in the total space
$\Sigma(x:A).\, B(x)$, viewed as a sheaf type via the projection to $A$.
\end{proposition}

\begin{proof}
The total space $\Sigma(x:A).B(x)$ is a sheaf if both $A$ and $B$ are
sheaves (sheaves are closed under $\Sigma$-types in a presheaf topos).
The descent data $(t_\alpha, b_\alpha)$ with patching $(p_{\alpha\beta},
q_{\alpha\beta})$ are exactly descent data for the total space.  Apply
Theorem~\ref{thm:descent-full}.
\end{proof}
